% =====================================================================
%  Guia de evaluacion: preguntas y respuestas
%  Proyecto: Sistema bidireccional Texto <-> Braille en FPGA (DE10-Lite)
%  Compilar:  pdflatex guia_evaluacion.tex   (dos veces)
% =====================================================================
\documentclass[11pt,a4paper]{article}

\usepackage[utf8]{inputenc}
\usepackage[T1]{fontenc}
\usepackage[spanish,es-noshorthands,es-tabla]{babel}
\usepackage{lmodern}
\usepackage{amsmath,amssymb}
\usepackage[margin=2.2cm]{geometry}
\usepackage{xcolor}
\usepackage{listings}
\usepackage{booktabs}
\usepackage{tabularx}
\usepackage{longtable}
\usepackage{array}
\usepackage{enumitem}
\usepackage[most]{tcolorbox}
\usepackage{fancyhdr}
\usepackage{titlesec}
\usepackage{tikz}
\usetikzlibrary{automata,positioning,arrows.meta,shapes.geometric,calc,fit,decorations.pathreplacing}
\usepackage[hidelinks]{hyperref}

% ---------------- Colores ----------------
\definecolor{tinta}{HTML}{13202B}
\definecolor{verde}{HTML}{137366}
\definecolor{verdesuave}{HTML}{DCEFEB}
\definecolor{ambar}{HTML}{9A520C}
\definecolor{ambarsuave}{HTML}{F8E6D3}
\definecolor{rojo}{HTML}{A12A2A}
\definecolor{rojosuave}{HTML}{F6DEDE}
\definecolor{azul}{HTML}{1F4E9A}
\definecolor{azulsuave}{HTML}{E1E9F6}
\definecolor{gris}{HTML}{5B6874}
\definecolor{fondocodigo}{HTML}{F4F6F7}

% ---------------- Titulos y encabezado ----------------
\titleformat{\section}{\Large\bfseries\color{tinta}}{\thesection}{0.8em}{}[\vspace{-0.4em}\textcolor{verde}{\rule{\linewidth}{1pt}}]
\titleformat{\subsection}{\large\bfseries\color{tinta}}{\thesubsection}{0.7em}{}
\pagestyle{fancy}
\fancyhf{}
\fancyhead[L]{\small\color{gris}Guía de evaluación \textendash{} Sistema Braille FPGA}
\fancyhead[R]{\small\color{gris}\leftmark}
\fancyfoot[C]{\small\thepage}

% ---------------- Codigo ----------------
\lstdefinestyle{vlog}{
  language=Verilog, basicstyle=\ttfamily\small,
  keywordstyle=\color{azul}\bfseries, commentstyle=\color{gris}\itshape,
  backgroundcolor=\color{fondocodigo}, frame=single, rulecolor=\color{gris!40},
  framesep=5pt, xleftmargin=6pt, xrightmargin=6pt,
  breaklines=true, columns=fullflexible, keepspaces=true, showstringspaces=false,
  literate={á}{{\'a}}1 {é}{{\'e}}1 {í}{{\'i}}1 {ó}{{\'o}}1 {ú}{{\'u}}1 {ñ}{{\~n}}1
}
\lstdefinestyle{mat}{style=vlog, language=Matlab}
\lstset{style=vlog}

% ---------------- Preguntas ----------------
\newcounter{preg}
\newcommand{\nivel}[1]{%
  \ifx#1B\colorbox{verdesuave}{\textcolor{verde}{\scriptsize\bfseries BÁSICA}}\fi
  \ifx#1M\colorbox{ambarsuave}{\textcolor{ambar}{\scriptsize\bfseries MEDIA}}\fi
  \ifx#1A\colorbox{rojosuave}{\textcolor{rojo}{\scriptsize\bfseries AVANZADA}}\fi}
\newcommand{\preg}[2]{\refstepcounter{preg}\par\vspace{0.9em}\noindent
  {\color{tinta}\bfseries P\thepreg.}~\nivel{#1}~\textbf{#2}\par\vspace{0.3em}}
\newtcolorbox{resp}{breakable,enhanced,colback=fondocodigo,colframe=verde,
  boxrule=0pt,leftrule=2.5pt,arc=0pt,left=8pt,right=6pt,top=4pt,bottom=4pt}
\newtcolorbox{clave}[1][Lo que el profesor quiere escuchar]{breakable,colback=verdesuave,colframe=verde,
  fonttitle=\bfseries\small,title=#1,arc=1.5mm,boxrule=0.6pt}
\newtcolorbox{trampa}[1][Cuidado: pregunta trampa]{breakable,colback=ambarsuave,colframe=ambar,
  fonttitle=\bfseries\small,title=#1,arc=1.5mm,boxrule=0.6pt}
\newcommand{\cod}[1]{\texttt{#1}}
\setlength{\emergencystretch}{3em}
\setlist{itemsep=1pt,topsep=3pt}
\renewcommand{\arraystretch}{1.2}

% ---------------- Estilos TikZ ----------------
\tikzset{
  >=Stealth,
  st/.style={circle,draw=verde,very thick,fill=verdesuave,minimum size=1.55cm,font=\small\ttfamily\bfseries,align=center},
  stt/.style={circle,draw=ambar,very thick,fill=ambarsuave,minimum size=1.55cm,font=\small\ttfamily\bfseries,align=center},
  bx/.style={rectangle,rounded corners=3pt,draw=gris,thick,fill=white,minimum height=0.9cm,minimum width=2.4cm,font=\small\ttfamily,align=center},
  bxr/.style={bx,draw=verde,fill=verdesuave},
  bxt/.style={bx,draw=ambar,fill=ambarsuave},
  bxc/.style={bx,draw=tinta,fill=tinta,text=white},
  sc/.style={rectangle,rounded corners=4pt,draw=azul,very thick,fill=azulsuave,minimum height=0.8cm,minimum width=1.9cm,font=\scriptsize\ttfamily\bfseries,align=center},
  lab/.style={font=\scriptsize,align=center},
  dec/.style={diamond,draw=ambar,thick,fill=ambarsuave,aspect=2.2,inner sep=1pt,font=\scriptsize,align=center},
  act/.style={rectangle,draw=verde,thick,fill=verdesuave,rounded corners=2pt,font=\scriptsize,align=center,inner sep=4pt}
}

% =====================================================================
\begin{document}

\begin{titlepage}
\centering
\vspace*{2.5cm}
{\color{verde}\large\bfseries PROYECTO FINAL \textendash{} DE10-LITE\par}
\vspace{0.6cm}
{\color{tinta}\Huge\bfseries Guía de evaluación\\[0.2cm]\LARGE Sistema bidireccional Texto $\leftrightarrow$ Braille\par}
\vspace{0.6cm}
{\large\color{gris}Banco de preguntas con respuestas, justificaciones y diagramas\par}
\vspace{1.4cm}
\begin{minipage}{0.82\textwidth}
\begin{clave}[Cómo usar esta guía]
Las preguntas están ordenadas por tema, del sistema completo a cada módulo, y
terminan con análisis crítico y ejercicios de traza. Cada una tiene un nivel:\\[3pt]
\nivel{B}\ entendimiento básico \quad \nivel{M}\ funcionamiento y flujo \quad \nivel{A}\ justificación, casos límite y análisis.\\[3pt]
Las cajas verdes resumen \textbf{lo que el evaluador espera escuchar}; las naranjas
señalan \textbf{preguntas trampa} o errores comunes. Practica respondiendo en voz
alta antes de leer la respuesta.
\end{clave}
\end{minipage}
\vfill
{\color{gris}Archivos evaluados: \cod{braille\_top.v}, \cod{braille\_controller.v}, \cod{uart\_rx.v}, \cod{uart\_tx.v},
\cod{line\_rx.v}, \cod{msg\_tx.v}, \cod{char\_buffer.sv}, \cod{button\_conditioner.v}, \cod{chord\_capture.v},
\cod{servo\_controller.v}, \cod{ascii\_to\_braille.v}, \cod{braille\_to\_ascii.v}, \cod{hex7seg.v}, \cod{matlab}\par}
\end{titlepage}

\tableofcontents
\newpage

% =====================================================================
\section{Visión general y arquitectura}

\preg{B}{Describe en un minuto qué hace tu proyecto.}
\begin{resp}
Es un sistema que traduce entre texto y Braille en las dos direcciones usando la
FPGA de la DE10-Lite. Tiene una celda Braille física de 6 servos, un teclado
Braille de 6 botones más un botón NEXT, y se comunica por UART a 9600 baudios con
una interfaz en MATLAB. Tiene tres modos útiles más reposo, elegidos con los
switches: \textbf{lectura} (la PC manda una palabra y la celda la muestra letra
por letra), \textbf{escritura} (el usuario teclea Braille y la PC recibe la
letra) y \textbf{práctica} (la FPGA elige letras al azar, compara la respuesta y
lleva aciertos y errores).
\end{resp}

\preg{B}{¿Dónde está la lógica del sistema: en la FPGA o en MATLAB? ¿Por qué?}
\begin{resp}
En la FPGA. MATLAB solo muestra, habla, grafica y envía tres comandos (\cod{Q},
\cod{R}, \cod{L}). La FPGA elige el modo (por switches), traduce, compara, cuenta,
genera las letras aleatorias y decide cuándo termina una sesión.
\end{resp}
\begin{clave}
Tres razones: (1) el objetivo del curso es el diseño de hardware; (2) la tarjeta
funciona sola, incluso sin PC, con los LEDs y displays; (3) la FPGA es la
\emph{única fuente de verdad}: el marcador y el modo nunca pueden quedar
desincronizados entre la PC y la tarjeta.
\end{clave}

\preg{M}{Dibuja la arquitectura y explica el camino de los datos.}
\begin{resp}
\begin{center}
\begin{tikzpicture}[node distance=0.5cm and 0.9cm]
\node[bx,minimum height=5.4cm,minimum width=1.6cm] (pc) {PC\\MATLAB};
\node[bxr,right=1.7cm of pc.north east,anchor=north west,yshift=-0.2cm] (rx) {uart\_rx};
\node[bxr,right=of rx] (lrx) {line\_rx};
\node[bx,below=of lrx] (buf) {char\_buffer};
\node[bx,below=0.5cm of rx] (btn) {button\_\\conditioner};
\node[bx,below=of buf] (ch) {chord\_capture};
\node[bxt,below=3.5cm of rx,minimum width=5.7cm,xshift=1.65cm] (tx) {msg\_tx + uart\_tx};
\node[bxc,right=1.1cm of lrx.north east,anchor=north west,minimum height=5.4cm,minimum width=2.3cm] (ct) {braille\_\\controller};
\node[bx,right=0.9cm of ct.north east,anchor=north west] (sv) {servo\_\\controller};
\node[bx,below=0.5cm of sv] (hx) {hex7seg\\LEDs};
\node[bx,below=0.5cm of hx] (sw) {switches\\(sync)};
\draw[->,thick,verde] (pc.east|-rx) -- node[above,lab]{UART\_RX} (rx);
\draw[->,thick,verde] (rx) -- (lrx);
\draw[->,thick] (lrx) -- (buf);
\draw[->,thick,verde] (lrx.east) -- (ct.west|-lrx);
\draw[->,thick] (buf.east) -- (ct.west|-buf);
\draw[->,thick] (ch.east) -- (ct.west|-ch);
\draw[->,thick] (btn.south) |- (ch.west);
\draw[->,thick,ambar] (ct.west|-tx) -- (tx.east);
\draw[->,thick,ambar] (tx.west) -- node[above,lab]{UART\_TX} (pc.east|-tx);
\draw[->,thick] (ct.east|-sv) -- (sv);
\draw[->,thick] (ct.east|-hx) -- (hx);
\draw[->,thick] (sw) -- (ct.east|-sw);
\end{tikzpicture}
\end{center}
\textbf{Entrada PC}: \cod{uart\_rx} (bits$\to$byte) $\to$ \cod{line\_rx}
(bytes$\to$línea; los datos van a \cod{char\_buffer}) $\to$ controller.
\textbf{Entrada física}: botones $\to$ \cod{button\_conditioner} $\to$
\cod{chord\_capture} $\to$ controller. \textbf{Salidas}: controller $\to$
\cod{servo\_controller} (celda), displays/LEDs y \cod{msg\_tx}$\to$\cod{uart\_tx}
(respuesta a la PC).
\end{resp}

\preg{M}{¿Cuántos módulos tiene el proyecto y cuál es la responsabilidad de cada uno en una frase?}
\begin{resp}
\begin{tabularx}{\linewidth}{>{\ttfamily}l X}
\toprule
\normalfont\textbf{Módulo} & \textbf{Responsabilidad} \\
\midrule
braille\_top & Conecta todo con los pines y fija los parámetros. \\
braille\_controller & Decide qué hace el sistema según modo y eventos. \\
uart\_rx / uart\_tx & Convierten bits del cable en un byte y viceversa. \\
line\_rx / msg\_tx & Convierten bytes en una línea de comando y un mensaje en bytes. \\
char\_buffer & Guarda la palabra de hasta 16 caracteres. \\
button\_conditioner & Sincroniza y quita el rebote de los 7 botones. \\
chord\_capture & Junta los botones presionados en un patrón (acorde). \\
servo\_controller & Genera el PWM de los 6 servos. \\
ascii\_to\_braille / braille\_to\_ascii & Tablas letra$\leftrightarrow$puntos. \\
hex7seg & Número de 4 bits $\to$ display de 7 segmentos. \\
\bottomrule
\end{tabularx}
Son 13 módulos de Verilog/SystemVerilog más el programa de MATLAB.
\end{resp}

\preg{M}{¿Qué criterio usaste para dividir el sistema en módulos?}
\begin{resp}
\textbf{Una responsabilidad por módulo} y separación por \textbf{niveles de
abstracción}. En la comunicación hay tres niveles: bits$\leftrightarrow$byte
(\cod{uart\_rx/tx}), bytes$\leftrightarrow$mensaje (\cod{line\_rx}/\cod{msg\_tx})
y significado del mensaje (controller). Así cada módulo se prueba por separado,
se puede cambiar uno sin tocar los demás (por ejemplo, la velocidad UART solo
afecta a \cod{CLKS\_PER\_BIT}) y el controller no sabe nada de tiempos de bit.
\end{resp}

\preg{B}{¿Qué es la convención ``bit $i$ = punto $i+1$'' y por qué es importante?}
\begin{resp}
En todos los módulos, el bit 0 es el punto 1, el bit 1 el punto 2, y así hasta el
bit 5 = punto 6. El botón \cod{DOT\_N[i]}, el servo \cod{SERVO[i]}, el LED
\cod{LEDR[i]} y el bit \cod{i} de los patrones son el mismo punto. Esto permite
pasar el patrón del teclado al controller y a los servos \textbf{sin reordenar
bits}, eliminando una fuente común de errores.
\end{resp}

\preg{B}{¿Qué hardware usa el proyecto y cómo se conecta?}
\begin{resp}
DE10-Lite (MAX 10, reloj de 50 MHz); 6 servos SG90 en \cod{SERVO[5:0]}; 6 botones
externos activos en bajo con pull-up en \cod{DOT\_N[5:0]}; KEY0 (reset) y KEY1
(NEXT); SW0--SW2 (modo y pista); un FT232RL (USB) o HC-05 (Bluetooth) en
\cod{UART\_RX}/\cod{UART\_TX} cruzados (el TX de uno al RX del otro); LEDs y
displays para depuración.
\end{resp}

% =====================================================================
\section{Fundamentos de diseño digital aplicados}

\preg{B}{¿Cuál es la diferencia entre lógica combinacional y secuencial? Da un ejemplo de cada una en tu proyecto.}
\begin{resp}
La \textbf{combinacional} depende solo de las entradas actuales, no tiene memoria
ni reloj: \cod{ascii\_to\_braille}, \cod{braille\_to\_ascii}, \cod{hex7seg} (todos
con \cod{always @(*)}), y la lectura de \cod{char\_buffer}. La
\textbf{secuencial} depende también del estado guardado y cambia en flancos de
reloj: \cod{uart\_rx}, los contadores de antirrebote, el contador PWM, la FSM del
controller (\cod{always @(posedge clk ...)}).
\end{resp}

\preg{M}{¿Por qué usas \texttt{=} en unos bloques y \texttt{<=} en otros?}
\begin{resp}
\cod{=} (bloqueante) en lógica combinacional \cod{always @(*)}: se evalúa en orden,
como un programa. \cod{<=} (no bloqueante) en lógica secuencial: todos los registros
toman su nuevo valor \emph{al mismo tiempo} al final del flanco, que es como se
comporta el hardware real. Mezclarlos en lógica secuencial puede causar
diferencias entre simulación y síntesis o condiciones de carrera.
\end{resp}
\begin{trampa}
Consecuencia práctica: en \cod{line\_rx}, \cod{buf\_waddr = len} usa el valor
\textbf{actual} de \cod{len}; el incremento \cod{len <= len + 1} se ve hasta el
siguiente ciclo. Por eso la primera letra se guarda en la casilla 0 y no en la 1.
\end{trampa}

\preg{M}{¿Qué es un latch inferido y cómo lo evitas?}
\begin{resp}
Si en un \cod{always @(*)} existe un camino donde una salida no se asigna, el
sintetizador debe ``recordar'' el valor anterior y crea un latch (memoria no
deseada, sensible a glitches y difícil de analizar en tiempos). Lo evito con un
\cod{default} en todos los \cod{case} y asignando valores por defecto al inicio
del bloque (por ejemplo \cod{braille\_out = 0; char\_valid = 1;} en
\cod{ascii\_to\_braille}).
\end{resp}

\preg{A}{¿Qué es la metaestabilidad y dónde la previenes?}
\begin{resp}
Cuando una señal asíncrona cambia justo dentro de la ventana de setup/hold de un
flip-flop, su salida puede quedar un tiempo en un valor intermedio antes de
resolverse a 0 o 1. Si ese valor indeciso se propaga, distintas partes del
circuito pueden leerlo diferente. Se previene con un \textbf{sincronizador de dos
flip-flops}: el primero puede quedar metaestable, pero tiene un ciclo completo
(20 ns) para resolverse antes de que lo muestree el segundo.
\begin{center}
\begin{tikzpicture}[node distance=1.4cm]
\node (in) {\small señal externa};
\node[bx,minimum width=1.4cm,right=of in] (f1) {FF1};
\node[bx,minimum width=1.4cm,right=of f1] (f2) {FF2};
\node[right=of f2] (out) {\small señal estable};
\draw[->,thick] (in) -- (f1);
\draw[->,thick] (f1) -- node[above,lab]{puede ser\\metaestable} (f2);
\draw[->,thick] (f2) -- (out);
\draw[->] ($(f1.south)+(0,-0.6)$) node[below,lab]{clk} -- (f1.south);
\draw[->] ($(f2.south)+(0,-0.6)$) node[below,lab]{clk} -- (f2.south);
\end{tikzpicture}
\end{center}
Está en: \cod{uart\_rx} (\cod{rx\_sync\_1/2}), \cod{button\_conditioner}
(\cod{sync\_ff1/2}) y \cod{braille\_top} para los switches (\cod{sw\_s1/2}).
\end{resp}

\preg{A}{¿Por qué \texttt{uart\_tx} no tiene sincronizador y \texttt{uart\_rx} sí?}
\begin{resp}
Porque la entrada de \cod{uart\_tx} (\cod{data\_in}, \cod{tx\_start}) viene de
\cod{msg\_tx}, que usa el mismo reloj de 50 MHz: ya es síncrona. En cambio, la
línea \cod{rx} viene de la PC, que tiene su propio reloj y no guarda ninguna
relación de fase con el de la FPGA.
\end{resp}

\preg{M}{¿Qué es un pulso de un ciclo y cómo se genera en tu código?}
\begin{resp}
Una señal que vale 1 durante exactamente un periodo de reloj (20 ns). Se usa para
avisar un evento \emph{una sola vez}. Patrón usado en todo el proyecto: asignar 0
como valor por defecto al inicio de cada ciclo y 1 solo en el caso que lo genera.
\begin{lstlisting}
data_ready <= 1'b0;                 // por defecto, cada ciclo
...
STOP: if (rx_sync_2) data_ready <= 1'b1;   // solo aqui
\end{lstlisting}
Ejemplos: \cod{data\_ready}, \cod{tx\_done}, \cod{press\_pulse}, \cod{chord\_valid},
\cod{line\_ack}, \cod{msg\_start}, \cod{uart\_start}.
\end{resp}

\preg{M}{¿Qué tipo de reset usas? ¿Es síncrono o asíncrono? ¿Activo en alto o en bajo?}
\begin{resp}
Asíncrono (\cod{always @(posedge clk or posedge rst)}) y activo en alto. El botón
KEY0 es activo en bajo, así que el top lo invierte: \cod{rst = \textasciitilde KEY[0]}.
Durante el reset todos los registros toman un valor conocido: FSMs en IDLE, líneas
TX en 1 (reposo UART), contadores en cero.
\end{resp}
\begin{trampa}[Pregunta de rigor]
\textbf{¿Tiene algún riesgo ese reset?} Sí: la \emph{liberación} de un reset
asíncrono también es asíncrona; si ocurre cerca de un flanco de reloj, algunos
flip-flops pueden salir del reset un ciclo antes que otros. La práctica
recomendada es un \emph{sincronizador de reset} (afirmación asíncrona, liberación
síncrona). En este proyecto no está implementado; es una mejora válida que conviene
mencionar.
\end{trampa}

\preg{B}{¿Qué es un \texttt{parameter} y para qué te sirvió?}
\begin{resp}
Una constante configurable de un módulo que el módulo padre puede sobrescribir al
instanciarlo con \cod{\#(.NOMBRE(valor))}. Todos los parámetros del sistema
(\cod{CLKS\_PER\_BIT}, \cod{DEBOUNCE\_COUNT}, \cod{PULSE\_*}, \cod{MAX\_LEN},
\cod{FEEDBACK\_DELAY}) se fijan en \cod{braille\_top}. Así se configura todo en
un solo lugar y se garantiza que los módulos que comparten un valor
(\cod{MAX\_LEN} en \cod{line\_rx}, \cod{char\_buffer} y el controller) usen el
mismo. También facilita simular con valores pequeños (por ejemplo, un antirrebote
de 10 ciclos en lugar de 500\,000).
\end{resp}
\begin{trampa}
Los valores por defecto dentro de un módulo \textbf{no se usan} si el top los
sobrescribe. Los anchos de pulso reales de los servos son los de \cod{braille\_top},
no los de \cod{servo\_controller}.
\end{trampa}

% =====================================================================
\section{Módulos básicos: tablas, displays y memoria}

\preg{B}{¿Qué es una LUT y cuáles hay en tu proyecto?}
\begin{resp}
Una \emph{Look-Up Table}: una tabla donde a cada entrada le corresponde una
salida fija. Se implementa con un \cod{case} combinacional. Hay tres:
\cod{ascii\_to\_braille} (letra $\to$ 6 puntos), \cod{braille\_to\_ascii}
(6 puntos $\to$ letra) y \cod{hex7seg} (4 bits $\to$ 7 segmentos).
\end{resp}

\preg{M}{¿Qué caracteres soporta \texttt{ascii\_to\_braille} y qué pasa con los demás?}
\begin{resp}
\textbf{a--z}, \textbf{A--Z} (mismo patrón que la minúscula, porque no se implementa
el signo de mayúscula de Braille), \textbf{ñ/Ñ} en Latin-1 (\cod{0xF1}/\cod{0xD1},
puntos 1-2-4-5-6) y el \textbf{espacio} (celda vacía, válido). Cualquier otro
carácter cae en \cod{default}: celda vacía y \cod{char\_valid = 0}. En modo lectura
la FPGA igual responde \cod{D} + ese carácter, pero no levanta ningún punto.
\end{resp}

\preg{A}{¿Por qué el patrón \texttt{000000} es válido en \texttt{ascii\_to\_braille} (espacio) pero inválido en \texttt{braille\_to\_ascii}?}
\begin{resp}
En la dirección letra$\to$puntos, el espacio se representa naturalmente como una
celda vacía. En la dirección puntos$\to$letra, \cod{chord\_capture} nunca entrega
\cod{000000} (solo entrega cuando hubo al menos un botón), y si fuera válido
``no presionar nada'' podría interpretarse como espacio. Por eso el espacio en
escritura se genera con el botón \textbf{NEXT}, de forma explícita.
\end{resp}

\preg{M}{¿Por qué la ñ debe llegar en Latin-1 y no en UTF-8?}
\begin{resp}
En UTF-8 la ñ ocupa dos bytes (\cod{0xC3 0xB1}); la FPGA procesa un byte por
carácter y vería dos caracteres desconocidos. En Latin-1 ocupa uno solo
(\cod{0xF1}), que es el que está en la tabla. MATLAB lo garantiza con
\cod{uint8(char(txt))}, que produce un byte por carácter (241 para la ñ).
\end{resp}

\preg{B}{¿Cómo funciona \texttt{hex7seg} y por qué los valores parecen invertidos?}
\begin{resp}
Es un \cod{case} de 16 entradas que devuelve \cod{seg = \{g,f,e,d,c,b,a\}}. Los
displays de la DE10-Lite son \textbf{activos en bajo}: un segmento se enciende con 0.
Por eso el 0 es \cod{1000000} (todo encendido menos \cod{g}). La entrada
\cod{blank} apaga el display completo (\cod{1111111}), usada en HEX4.
\end{resp}

\preg{M}{¿Por qué los aciertos se ven en hexadecimal en los displays?}
\begin{resp}
\cod{hits} es de 8 bits y se parte en dos nibbles: \cod{hits[7:4]} a HEX3 y
\cod{hits[3:0]} a HEX2. Cada \cod{hex7seg} muestra 0--F, así que 26 aciertos se
ven como \cod{1A}. Mostrarlo en decimal requeriría convertir binario a BCD
(divisiones entre 10 o el algoritmo \emph{double dabble}); se priorizó la
simplicidad.
\end{resp}

\preg{M}{Explica \texttt{char\_buffer}: tamaño, puertos y tipo de lectura.}
\begin{resp}
Memoria de 16 casillas de 8 bits (\cod{reg [7:0] mem [0:15]}, 128 bits) con un
puerto de escritura y otro de lectura independientes. \textbf{Escritura síncrona}
(\cod{always @(posedge clk) if (wr\_en) mem[wr\_addr] <= wr\_data}), hecha por
\cod{line\_rx}. \textbf{Lectura combinacional} (\cod{assign rd\_data = mem[rd\_addr]}),
hecha por el controller: el dato está disponible en el mismo ciclo en que se pone
la dirección.
\end{resp}

\preg{A}{¿Qué implicación tiene la lectura combinacional en la síntesis? ¿Por qué no tiene reset?}
\begin{resp}
Los bloques de memoria M9K de la MAX 10 solo permiten lectura \textbf{síncrona}
(registrada). Con lectura combinacional, Quartus no puede usarlos y construye la
memoria con flip-flops y multiplexores en lógica. Para 128 bits el costo es
despreciable y se gana simplicidad (no hay un ciclo de latencia que manejar en la
FSM). Para memorias de miles de bytes convendría lectura síncrona.\par
No tiene reset porque no lo necesita: el controller lleva \cod{word\_len}, la
cantidad de caracteres válidos, y nunca lee más allá. Además, las memorias con
reset no se pueden mapear a bloques de RAM.
\end{resp}

\preg{B}{¿Por qué \texttt{AW = 4}?}
\begin{resp}
\emph{Address Width}: bits de dirección. $2^4 = 16$ casillas = \cod{MAX\_LEN}. Para
32 caracteres serían \cod{MAX\_LEN = 32, AW = 5}, cambiando solo el top (y el
\cod{MAX\_LEN} de MATLAB).
\end{resp}

% =====================================================================
\section{UART: capa física}

\preg{B}{¿Qué es UART y qué significa 9600 8N1?}
\begin{resp}
\emph{Universal Asynchronous Receiver-Transmitter}: comunicación serial sin línea
de reloj; ambos lados acuerdan la velocidad. 9600 = bits por segundo (baudios);
8N1 = 8 bits de datos, sin paridad (\emph{No parity}), 1 bit de stop.
\end{resp}

\preg{M}{Dibuja la trama de un byte. ¿Por qué hay bit de START y de STOP?}
\begin{resp}
Ejemplo con \cod{'R'} = \cod{0x52} = \cod{0101\,0010}, enviado LSB primero:
\begin{center}
\begin{tikzpicture}[xscale=1.05,yscale=0.9]
\foreach \i/\l in {0/reposo,1/START,2/D0,3/D1,4/D2,5/D3,6/D4,7/D5,8/D6,9/D7,10/STOP,11/reposo}{
  \draw[gris!40] (\i,0) -- (\i,1.4);
  \node[lab] at (\i+0.5,-0.35) {\l};}
\draw[gris!40] (12,0) -- (12,1.4);
\foreach \i/\b in {2/0,3/1,4/0,5/0,6/1,7/0,8/1,9/0}{\node[font=\scriptsize\ttfamily] at (\i+0.5,1.65) {\b};}
\draw[very thick,verde] (0,1) -- (1,1) -- (1,0) -- (3,0) -- (3,1) -- (4,1) -- (4,0) -- (6,0) -- (6,1) -- (7,1) -- (7,0) -- (8,0) -- (8,1) -- (9,1) -- (9,0) -- (10,0) -- (10,1) -- (12,1);
\foreach \i in {1,...,10}{\fill[ambar] (\i+0.5,0.5) circle (1.6pt);}
\draw[decorate,decoration={brace,mirror,amplitude=4pt}] (2,-0.7) -- (3,-0.7) node[midway,below=4pt,lab]{1 bit = 5208 ciclos = 104\,µs};
\draw[decorate,decoration={brace,mirror,amplitude=4pt}] (1,-1.4) -- (11,-1.4) node[midway,below=4pt,lab]{10 bits = 1.04 ms};
\end{tikzpicture}
\end{center}
La línea descansa en 1. El \textbf{START} (0) crea un flanco que marca el inicio
del byte y sincroniza al receptor (no hay reloj compartido). El \textbf{STOP} (1)
deja la línea en reposo para que el siguiente START sea detectable, y permite
detectar errores de trama.
\end{resp}

\preg{B}{¿De dónde sale 5208 y por qué el contador es de 13 bits?}
\begin{resp}
$50\,000\,000 / 9600 = 5208.33 \to 5208$ ciclos por bit. El error es
$0.33/5208.33 \approx 0.006\%$, muy por debajo de la tolerancia de UART
(del orden de $\pm 2$--$5\%$ acumulado en una trama de 10 bits). El contador
debe llegar a 5207: $2^{12} = 4096$ no alcanza, $2^{13} = 8192$ sí.
\end{resp}

\preg{M}{Dibuja la FSM de \texttt{uart\_rx} y explica cada transición.}
\begin{resp}
\begin{center}
\begin{tikzpicture}[->,node distance=2.5cm,auto]
\node[st,initial,initial text=] (I) {IDLE};
\node[st,right=of I] (S) {START};
\node[st,right=of S] (D) {DATA};
\node[st,right=of D] (P) {STOP};
\path (I) edge node[lab,above]{rx = 0} (S)
      (S) edge node[lab,above]{½ bit\\rx = 0} (D)
      (S) edge[bend left=45] node[lab,below]{½ bit, rx = 1 (ruido)} (I)
      (D) edge[loop above] node[lab]{cada 5208 ciclos\\guarda bit; bit\_index++} (D)
      (D) edge node[lab,above]{bit 7} (P)
      (P) edge[bend right=-60] node[lab,below]{1 bit: si rx = 1 $\Rightarrow$ data\_out, data\_ready} (I);
\end{tikzpicture}
\end{center}
\textbf{IDLE}: espera que la línea sincronizada baje a 0. \textbf{START}: espera
medio bit (2604 ciclos) y confirma el 0; si ya es 1, era ruido y regresa.
\textbf{DATA}: cada 5208 ciclos guarda \cod{rx\_data[bit\_index] <= rx\_sync\_2};
tras el bit 7 pasa a STOP. \textbf{STOP}: un bit después, si la línea vale 1,
copia el byte a \cod{data\_out} y da el pulso \cod{data\_ready}; si vale 0 (error de
trama), descarta el byte.
\end{resp}

\preg{A}{¿Por qué se espera medio bit en START? ¿Qué pasaría si se esperara uno completo?}
\begin{resp}
Para muestrear en el \textbf{centro} de cada bit, donde la señal es más estable y
lejos de las transiciones. Después de medio bit se está en el centro del START;
sumando 5208 ciclos se cae en el centro de D0, D1, etc. Además sirve como filtro
de ruido: un pico corto no sobrevive medio bit.\par
Si se esperara un bit completo, cada muestra caería justo en el borde entre bits:
cualquier pequeña diferencia de reloj entre PC y FPGA, o el retraso del
sincronizador, haría leer el bit vecino. El margen de error sería casi nulo.
\end{resp}

\preg{A}{¿Cuánto error de reloj tolera tu receptor?}
\begin{resp}
El receptor se resincroniza en cada START y la última muestra es la del STOP, a
$9.5$ bits del flanco de inicio. Para no salirse del bit, el desplazamiento
acumulado debe ser menor a medio bit: $9.5 \cdot \varepsilon < 0.5 \Rightarrow
\varepsilon < 5.3\%$ (en la práctica se considera $\pm 2$--$3\%$ por cada lado,
porque los dos relojes contribuyen). El error de redondeo de 5208 es $0.006\%$,
y el sincronizador de 2 ciclos (40 ns) es despreciable frente a 104\,µs.
\end{resp}

\preg{M}{¿Por qué los bits se guardan con \texttt{rx\_data[bit\_index]} empezando en 0?}
\begin{resp}
Porque la UART transmite primero el bit menos significativo (LSB). El primer bit
que llega es el bit 0 del byte. \cod{uart\_tx} hace lo mismo con
\cod{tx <= tx\_data[bit\_index]}, así ambos lados coinciden.
\end{resp}

\preg{M}{¿Por qué existe \texttt{rx\_data} si ya existe \texttt{data\_out}?}
\begin{resp}
\cod{rx\_data} es el registro donde se \emph{arma} el byte bit por bit.
\cod{data\_out} solo se actualiza cuando el byte está completo \emph{y} el STOP es
válido. Así la salida nunca muestra un byte a medias ni uno con error de trama.
\end{resp}

\preg{M}{Dibuja la FSM de \texttt{uart\_tx}. ¿Por qué copia el byte en \texttt{tx\_data}?}
\begin{resp}
\begin{center}
\begin{tikzpicture}[->,node distance=2.5cm,auto]
\node[stt,initial,initial text=] (I) {IDLE\\[-2pt]\tiny tx=1};
\node[stt,right=of I] (S) {START\\[-2pt]\tiny tx=0};
\node[stt,right=of S] (D) {DATA\\[-2pt]\tiny tx=bit};
\node[stt,right=of D] (P) {STOP\\[-2pt]\tiny tx=1};
\path (I) edge node[lab,above]{tx\_start\\tx\_data$\leftarrow$data\_in} (S)
      (S) edge node[lab,above]{5208\\ciclos} (D)
      (D) edge[loop above] node[lab]{8 bits × 5208} (D)
      (D) edge node[lab,above]{bit 7} (P)
      (P) edge[bend right=-60] node[lab,below]{5208 ciclos: tx\_done (pulso), tx\_busy = 0} (I);
\end{tikzpicture}
\end{center}
Copia el byte porque transmitirlo tarda 1.04 ms; si \cod{data\_in} cambiara en ese
tiempo, el byte saldría mezclado. Con la copia, la entrada puede cambiar sin
afectar la transmisión en curso.
\end{resp}

\preg{A}{Compara \texttt{uart\_rx} y \texttt{uart\_tx}. ¿Por qué START dura distinto en cada uno?}
\begin{resp}
\begin{tabularx}{\linewidth}{l X X}
\toprule
& \cod{uart\_rx} & \cod{uart\_tx} \\
\midrule
Dirección & cable $\to$ byte & byte $\to$ cable \\
Sincronizador & sí (entrada externa) & no (entrada síncrona) \\
START & ½ bit (posicionarse al centro) & 1 bit (sostener el nivel) \\
Pulso de aviso & \cod{data\_ready} & \cod{tx\_done} \\
Verificación & revisa el STOP & --- \\
\bottomrule
\end{tabularx}
El receptor necesita \emph{alinearse} con los centros; el transmisor necesita
\emph{generar} bits de duración exacta.
\end{resp}

\preg{M}{¿Qué pasa si llega un byte con el STOP en 0?}
\begin{resp}
Es un error de trama (baudios distintos, ruido, cable mal conectado). El receptor
no levanta \cod{data\_ready}: el byte se descarta y vuelve a IDLE. Se prefiere perder
un byte a entregar uno incorrecto. No hay notificación de error hacia arriba
(sería una mejora).
\end{resp}

% =====================================================================
\section{UART: enmarcado y protocolo}

\preg{B}{¿La UART envía paquetes? ¿Cuántos bytes manda a la vez?}
\begin{resp}
La UART manda \textbf{un byte por trama}; no conoce paquetes. Los ``paquetes'' son
una convención del protocolo propio: una \textbf{línea} de texto que termina en
\cod{\textbackslash n} (byte 10). \cod{Rhola\textbackslash n} son 6 tramas seguidas.
\end{resp}

\preg{M}{Explica el formato del protocolo y por qué se eligió así.}
\begin{resp}
Cada mensaje es texto ASCII: \textbf{una letra de comando + datos opcionales +
\cod{\textbackslash n}}. Ventajas: largo variable sin tener que enviar el largo
primero; legible en cualquier monitor serial (depuración fácil); el
\cod{\textbackslash n} es un delimitador inequívoco; MATLAB lo maneja nativamente
con \cod{configureTerminator(s,"LF")}.
\end{resp}

\preg{B}{Lista los comandos PC$\to$FPGA y las respuestas.}
\begin{resp}
\begin{tabularx}{\linewidth}{>{\ttfamily}l X >{\ttfamily}l l}
\toprule
\normalfont\textbf{Cmd} & \textbf{Significado} & \normalfont\textbf{Resp.} & \textbf{¿MATLAB?} \\
\midrule
R<texto> & Read: cargar y mostrar palabra (solo modo lectura) & D<c> & Sí \\
L<n> & Letras por sesión de práctica, 0--99 (0 = sin límite) & L<hh> & Sí \\
Q & Query: modo actual & M<m> & Sí \\
S & Score: marcador & S<hh><mm> & Pruebas \\
C & Clear: borrar marcador & S0000 & Pruebas \\
N & Next: igual que el botón & según modo & Pruebas \\
\bottomrule
\end{tabularx}
Errores: \cod{!L} (más de 16 caracteres), \cod{!U} (comando desconocido o mal
formado), \cod{!M} (\cod{R} fuera de modo lectura).
\end{resp}

\preg{B}{Lista los mensajes FPGA$\to$PC.}
\begin{resp}
\cod{M<m>} modo (I, R, W, P); \cod{D<c>} letra mostrada; \cod{E} fin de palabra;
\cod{K<c>} letra tecleada (\cod{?} inválida, espacio = NEXT); \cod{T<c>} letra
objetivo; \cod{O<c><hh><mm>} acierto; \cod{X<c><hh><mm>} error; \cod{F<hh><mm>}
fin de sesión; \cod{L<hh>} confirmación; \cod{S<hh><mm>} marcador; \cod{!L !U !M}
errores. Todos terminan en \cod{\textbackslash n} y miden de 2 a 7 bytes.
\end{resp}

\preg{A}{¿Por qué los contadores se envían como dos caracteres hexadecimales y no como un byte?}
\begin{resp}
Si se enviara el valor crudo, \cod{hits = 10} produciría el byte \cod{0x0A}, que es
\cod{\textbackslash n}: el receptor cortaría el mensaje. Igual con 13
(\cod{\textbackslash r}) o 0. En hexadecimal cada dígito es \cod{'0'--'9'}
(48--57) o \cod{'A'--'F'} (65--70); ninguno coincide con un delimitador.
Además: longitud fija (siempre 2 caracteres), legible, y en hardware es casi
gratis (partir en nibbles y sumar un desplazamiento); en decimal habría que
dividir entre 10.
\begin{lstlisting}
function [7:0] hexc; input [3:0] n;
  hexc = (n < 10) ? (8'h30 + n) : (8'h37 + n);   // '0'..'9' | 'A'..'F'
endfunction
\end{lstlisting}
Costo: 2 bytes extra por contador ($\approx$2 ms por mensaje), irrelevante frente al
ritmo humano.
\end{resp}

\preg{M}{Explica el flujo de \texttt{line\_rx}.}
\begin{resp}
\begin{center}
\begin{tikzpicture}[node distance=0.55cm and 0.9cm,->]
\node[act] (b) {llega byte\\(rx\_valid)};
\node[dec,below=of b] (lr) {¿line\_ready\\= 1?};
\node[act,right=of lr] (ig) {ignorar};
\node[dec,below=of lr] (t) {¿\textbackslash n o \textbackslash r?};
\node[dec,right=of t] (hc) {¿have\_cmd?};
\node[act,right=of hc] (ok) {line\_ready = 1};
\node[dec,below=of t] (fc) {¿have\_cmd\\= 0?};
\node[act,right=of fc] (cm) {cmd = byte\\len = 0, ovf = 0};
\node[dec,below=of fc] (le) {¿len < 16?};
\node[act,right=of le] (da) {len + 1\\(R/L: al buffer)};
\node[act,below=of le] (of) {overflow = 1};
\draw (b) -- (lr); \draw (lr) -- node[lab,above]{sí} (ig);
\draw (lr) -- node[lab,left]{no} (t); \draw (t) -- node[lab,above]{sí} (hc);
\draw (hc) -- node[lab,above]{sí} (ok);
\draw (t) -- node[lab,left]{no} (fc); \draw (fc) -- node[lab,above]{sí} (cm);
\draw (fc) -- node[lab,left]{no} (le); \draw (le) -- node[lab,above]{sí} (da);
\draw (le) -- node[lab,left]{no} (of);
\end{tikzpicture}
\end{center}
El primer byte de cada línea es el comando; los siguientes son datos (solo los de
\cod{R} y \cod{L} se escriben en el buffer, en la casilla \cod{len}); el
terminador levanta \cod{line\_ready} si hubo comando (las líneas vacías, como el
\cod{\textbackslash n} de un \cod{\textbackslash r\textbackslash n}, se ignoran). Más
de 16 datos marcan \cod{overflow}.
\end{resp}

\preg{A}{¿Por qué \texttt{len} tiene 5 bits si el buffer tiene direcciones de 4?}
\begin{resp}
\cod{len} cuenta cuántos datos llegaron y debe poder valer 16 (palabra llena),
que en binario es \cod{10000}: 5 bits (\cod{[AW:0]}). La dirección del buffer usa
solo los 4 bits bajos (\cod{len[3:0]}), porque cuando \cod{len = 16} ya no se escribe
(\cod{len < MAX\_LEN} es falso).
\end{resp}

\preg{A}{Escribe y explica la condición de \texttt{buf\_we}.}
\begin{resp}
\begin{lstlisting}
assign buf_we = accept && have_cmd && !is_term && (len < MAX_LEN) && store_cmd;
\end{lstlisting}
Cinco condiciones: el byte es nuevo y no hay línea pendiente (\cod{accept});
ya pasó el comando (es un dato); no es terminador; cabe; y el comando es \cod{R} o
\cod{L}. Es combinacional para que el byte se escriba en el mismo ciclo en que
llega, usando el \cod{len} actual como dirección.
\end{resp}

\preg{A}{Explica el handshake \texttt{line\_ready}/\texttt{ack}. ¿Por qué \texttt{line\_ready} es un nivel y no un pulso?}
\begin{resp}
\begin{center}
\begin{tikzpicture}[yscale=0.55,xscale=0.85]
\node[lab,anchor=east] at (-0.2,4.5) {rx\_valid};
\node[lab,anchor=east] at (-0.2,2.5) {line\_ready};
\node[lab,anchor=east] at (-0.2,0.5) {ack};
\draw[thick,verde] (0,4) -- (1,4) -- (1,5) -- (1.2,5) -- (1.2,4) -- (2.5,4) -- (2.5,5) -- (2.7,5) -- (2.7,4) -- (4,4) -- (4,5) -- (4.2,5) -- (4.2,4) -- (12,4);
\node[lab] at (4.1,5.6) {\textbackslash n};
\draw[thick,azul] (0,2) -- (4.4,2) -- (4.4,3) -- (8,3) -- (8,2) -- (12,2);
\draw[thick,ambar] (0,0) -- (7.8,0) -- (7.8,1) -- (8,1) -- (8,0) -- (12,0);
\draw[dashed,gris] (4.4,-0.4) -- (4.4,5.8);
\draw[dashed,gris] (7.8,-0.4) -- (7.8,5.8);
\node[lab] at (6.1,3.6) {esperando al controller};
\node[lab] at (10,1.6) {acepta otra línea};
\end{tikzpicture}
\end{center}
Al terminar una línea, \cod{line\_rx} levanta \cod{line\_ready} y \textbf{la
mantiene} hasta que el controller, en ST\_IDLE, copia \cod{cmd/len/ovf} y responde con
el pulso \cod{ack}. Si fuera un pulso de un ciclo y el controller estuviera ocupado
(por ejemplo en ST\_SEND), la línea se perdería. Mientras \cod{line\_ready = 1}, los
bytes nuevos se descartan (\cod{accept = rx\_valid \&\& !line\_ready}); por eso la PC
espera la respuesta antes de enviar otro comando.
\end{resp}

\preg{M}{Dibuja la FSM de \texttt{msg\_tx} y explica \texttt{data\_reg[idx*8 +: 8]}.}
\begin{resp}
\begin{center}
\begin{tikzpicture}[->,node distance=3cm,auto]
\node[stt,initial,initial text=] (I) {IDLE};
\node[stt,right=of I] (S) {SEND};
\node[stt,right=of S] (W) {WAIT};
\path (I) edge node[lab,above]{start: copia mensaje\\idx = 0, busy = 1} (S)
      (S) edge[bend left=20] node[lab,above]{uart\_data, uart\_start} (W)
      (W) edge[bend left=20] node[lab,below]{uart\_done y idx < len\\idx++} (S)
      (W) edge[loop right] node[lab]{espera} (W)
      (W) edge[bend left=55] node[lab,below]{uart\_done y idx = len: busy = 0} (I);
\end{tikzpicture}
\end{center}
\cod{data\_reg[idx*8 +: 8]} es un \emph{indexed part-select}: ``8 bits empezando en
el bit \cod{idx·8}''. Con \cod{idx = 0} toma \cod{[7:0]}, con 1 toma \cod{[15:8]},
etc. Cuando \cod{idx == len\_reg} ya no hay caracteres y envía \cod{8'h0A}
(\cod{\textbackslash n}). Por eso con \cod{msg\_len = 2} salen 3 bytes.
\end{resp}

\preg{A}{En \texttt{msg\_data <= \{32'd0, buf\_rdata, "D"\}}, ¿por qué la \texttt{"D"} va al final si se envía primero?}
\begin{resp}
En una concatenación de Verilog el último elemento ocupa los bits menos
significativos. \cod{"D"} queda en \cod{msg\_data[7:0]}, que es el byte 0, el primero
que envía \cod{msg\_tx}. La letra queda en \cod{[15:8]} (segundo byte). Los 32 ceros
rellenan hasta 48 bits y no se envían porque \cod{msg\_len = 2}.
\end{resp}

\preg{A}{¿Por qué \texttt{msg\_tx} copia el mensaje? ¿Qué ganas?}
\begin{resp}
Un mensaje tarda hasta 7.3 ms ($\approx$365\,000 ciclos). Con la copia en
\cod{data\_reg}, el controller sale de ST\_SEND en 2 ciclos y queda libre para
atender botones o preparar otra respuesta en \cod{msg\_data} sin corromper la que
se está enviando. Desacopla un proceso de nanosegundos (controller) de uno de
milisegundos (cable).
\end{resp}

\preg{M}{¿Por qué existen \texttt{ST\_SEND} y \texttt{ST\_GAP} en el controller?}
\begin{resp}
\textbf{ST\_SEND} espera \cod{msg\_busy = 0} antes de dar \cod{msg\_start}; si
\cod{msg\_tx} está ocupado, se queda ahí. \textbf{ST\_GAP} agrega un ciclo porque
\cod{msg\_busy} es un registro: sube un ciclo después de \cod{start}. Sin GAP, el
controller podría volver a IDLE, preparar otro mensaje y ver todavía
\cod{msg\_busy = 0}, disparando un segundo \cod{start} que \cod{msg\_tx} ignoraría
(está en SEND), y ese mensaje se perdería.
\end{resp}

\preg{M}{¿Cuántos bytes guarda y cuántos maneja cada módulo de la UART?}
\begin{resp}
\begin{tabularx}{\linewidth}{>{\ttfamily}l X X}
\toprule
\normalfont\textbf{Módulo} & \textbf{Guarda} & \textbf{Maneja por mensaje} \\
\midrule
uart\_rx & 1 byte (más el que arma) & 1 byte por trama \\
line\_rx & 1 byte (\cod{cmd}) + contador & 2 a 18 bytes por línea \\
char\_buffer & 16 bytes & datos de R/L \\
controller & \cod{cmd\_r} (1) + \cod{msg\_data} (6) & --- \\
msg\_tx & \cod{data\_reg} (6) + \cod{uart\_data} (1) & 2 a 7 bytes \\
uart\_tx & 1 byte & 1 byte por trama \\
\bottomrule
\end{tabularx}
\end{resp}

\preg{A}{¿Qué pasa si la PC envía dos comandos seguidos sin esperar respuesta?}
\begin{resp}
Depende del tiempo. El controller da \cod{ack} en el primer ciclo en ST\_IDLE, unos
pocos ciclos después del \cod{\textbackslash n}; un byte nuevo tarda 1.04 ms en
llegar, así que en la práctica el segundo comando casi siempre se acepta. El riesgo
real está si el controller estuviera ocupado en ST\_SEND esperando a que
\cod{msg\_tx} termine un mensaje largo: \cod{line\_ready} seguiría arriba y los
bytes del segundo comando que lleguen en ese tiempo se descartarían, produciendo una
línea mal formada. El protocolo, y MATLAB, lo evitan esperando la respuesta.
\end{resp}

\preg{M}{Traza completa: ¿qué pasa cuando MATLAB envía \texttt{Rhola} en modo lectura?}
\begin{resp}
\begin{center}
\begin{tikzpicture}[yscale=0.62]
\node[bx,minimum width=2cm] (p) at (0,0) {MATLAB};
\node[bxc,minimum width=2cm] (f) at (8.2,0) {FPGA};
\draw[gris,dashed] (0,-0.6) -- (0,-9.6);
\draw[gris,dashed] (8.2,-0.6) -- (8.2,-9.6);
\draw[->,thick,verde] (0,-1.4) -- node[above,lab]{\ttfamily Rhola\textbackslash n (6 bytes, 6.25 ms)} (8.2,-1.4);
\node[lab,anchor=west] at (8.4,-2.2) {uart\_rx $\to$ line\_rx: cmd=R,\\h,o,l,a $\to$ casillas 0--3};
\node[lab,anchor=west] at (8.4,-3.6) {controller: ack, word\_len=4,\\ST\_SHOW: 'h'$\to$010011 (servos)};
\draw[<-,thick,ambar] (0,-4.8) -- node[above,lab]{\ttfamily Dh\textbackslash n} (8.2,-4.8);
\node[lab,anchor=east] at (-0.2,-5.4) {muestra H y la dice};
\node[lab,anchor=west] at (8.4,-6.4) {NEXT: char\_idx = 1};
\draw[<-,thick,ambar] (0,-6.9) -- node[above,lab]{\ttfamily Do\textbackslash n} (8.2,-6.9);
\node[lab,anchor=west] at (8.4,-8.0) {... NEXT tras la última};
\draw[<-,thick,ambar] (0,-8.8) -- node[above,lab]{\ttfamily E\textbackslash n (servos abajo)} (8.2,-8.8);
\end{tikzpicture}
\end{center}
\end{resp}

% =====================================================================
\section{Entrada física: botones y acordes}

\preg{B}{¿Qué tres problemas resuelve \texttt{button\_conditioner}?}
\begin{resp}
(1) \textbf{Rebote}: los contactos mecánicos vibran unos milisegundos y una
presión parecería varias. (2) \textbf{Asincronía}: el botón no sigue el reloj
(riesgo de metaestabilidad). (3) \textbf{Lógica activa en bajo}: los botones dan 0
al presionar; el módulo invierte para que 1 = presionado.
\end{resp}

\preg{A}{Explica el algoritmo de antirrebote con un diagrama de tiempo.}
\begin{resp}
\begin{center}
\begin{tikzpicture}[yscale=0.55,xscale=0.9]
\node[lab,anchor=east] at (-0.2,4.5) {botón (crudo)};
\node[lab,anchor=east] at (-0.2,2.5) {contador};
\node[lab,anchor=east] at (-0.2,0.5) {pressed};
\draw[thick,gris] (0,5) -- (1,5) -- (1,4) -- (1.3,4) -- (1.3,5) -- (1.5,5) -- (1.5,4) -- (1.7,4) -- (1.7,5) -- (1.9,5) -- (1.9,4) -- (12,4);
\draw[thick,azul] (0,2) -- (1,2) -- (1.3,2.3) -- (1.3,2) -- (1.5,2) -- (1.7,2.3) -- (1.7,2) -- (1.9,2) -- (7,3) -- (7,2) -- (12,2);
\draw[thick,verde] (0,0) -- (7,0) -- (7,1) -- (12,1);
\draw[dashed,gris] (1.9,-0.4) -- (1.9,5.6);
\draw[dashed,gris] (7,-0.4) -- (7,5.6);
\draw[<->] (1.9,5.6) -- node[above,lab]{estable 10 ms = 500\,000 ciclos} (7,5.6);
\node[lab] at (1.45,3.2) {rebotes:\\se reinicia};
\end{tikzpicture}
\end{center}
Cada botón tiene su contador. Si el valor sincronizado coincide con el aceptado
(\cod{buttons\_pressed}), el contador se reinicia. Si difiere, cuenta; solo cuando
llega a 500\,000 ciclos seguidos (10 ms) se acepta el cambio. Un rebote devuelve el
valor al aceptado y reinicia el conteo. Al aceptar una presión se genera
\cod{press\_pulse} de un ciclo.
\end{resp}

\preg{M}{¿Por qué 10 ms y por qué el contador es de 19 bits?}
\begin{resp}
Los rebotes típicos de botones mecánicos duran menos de 5--10 ms; 10 ms los cubre
y sigue siendo imperceptible para una persona. $10\text{ ms} \times 50\text{ MHz}
= 500\,000$; $2^{19} = 524\,288 \geq 500\,000$ (con 18 bits, 262\,144, no alcanza).
\end{resp}

\preg{M}{¿Por qué \texttt{button\_conditioner} entrega dos salidas distintas?}
\begin{resp}
\cod{buttons\_pressed} (nivel) indica qué botones están presionados \emph{ahora}:
lo necesita \cod{chord\_capture} para saber cuáles se presionan juntos.
\cod{press\_pulse} (pulso) marca el instante de la presión: lo necesita NEXT para
avanzar \emph{una sola vez} por presión. Si NEXT usara el nivel, mantenerlo
presionado avanzaría en cada ciclo de reloj.
\end{resp}

\preg{M}{¿Cómo sabe la FPGA qué botones se presionan al mismo tiempo?}
\begin{resp}
Cada botón es un bit de \cod{buttons\_pressed[5:0]}; el vector completo muestra los
seis a la vez. El problema es que los dedos nunca presionan todo en el mismo
instante; eso lo resuelve \cod{chord\_capture}.
\end{resp}

\preg{A}{Explica \texttt{chord\_capture} con un diagrama de tiempo. ¿Por qué OR y por qué entregar al soltar?}
\begin{resp}
\begin{center}
\begin{tikzpicture}[yscale=0.5,xscale=0.9]
\foreach \y/\l in {6.5/punto 1,4.5/punto 2,2.5/punto 5,0.5/chord\_valid}{\node[lab,anchor=east] at (-0.2,\y) {\l};}
\draw[thick,verde] (0,6) -- (1,6) -- (1,7) -- (6,7) -- (6,6) -- (12,6);
\draw[thick,verde] (0,4) -- (2,4) -- (2,5) -- (7.5,5) -- (7.5,4) -- (12,4);
\draw[thick,verde] (0,2) -- (3,2) -- (3,3) -- (9,3) -- (9,2) -- (12,2);
\draw[thick,ambar] (0,0) -- (9.1,0) -- (9.1,1) -- (9.4,1) -- (9.4,0) -- (12,0);
\node[lab] at (4.5,8) {acc = 000001 $\to$ 000011 $\to$ 010011 (se conserva al soltar)};
\node[lab] at (10.8,1.5) {chord = 010011 (h)};
\end{tikzpicture}
\end{center}
\begin{lstlisting}
if (dots_pressed != 0) acc <= acc | dots_pressed;      // acumula
else if (acc != 0) begin chord <= acc; chord_valid <= 1; acc <= 0; end
\end{lstlisting}
El \textbf{OR} conserva cualquier punto que se haya presionado aunque sea un instante,
aunque se suelte antes que otros. Se entrega \textbf{al soltar todos} porque es el
único momento en que se sabe que el usuario terminó la letra. Como la entrada ya
viene sin rebote, la liberación es limpia y no genera acordes falsos.
\end{resp}

\preg{A}{¿Qué limitaciones tiene ese método?}
\begin{resp}
Si el usuario presiona un punto por error y lo suelta, queda incluido en el acorde
(no hay forma de ``borrar'' un punto a mitad de la letra). Si nunca se sueltan todos,
no se entrega nada. Son compromisos aceptables para un teclado Braille de
estilo Perkins, donde la letra se forma presionando simultáneamente.
\end{resp}

% =====================================================================
\section{Salida física: servos y PWM}

\preg{B}{¿Cómo se controla un servo SG90?}
\begin{resp}
Con PWM de 50 Hz (periodo de 20 ms). La duración del pulso en alto, entre $\sim$1 y
2 ms, define el ángulo; 1.5 ms es el centro.
\begin{center}
\begin{tikzpicture}[yscale=0.6,xscale=0.55]
\node[lab,anchor=east] at (-0.2,2.5) {punto arriba};
\node[lab,anchor=east] at (-0.2,0.5) {punto abajo};
\draw[thick,verde] (0,2) -- (0,3) -- (1.16,3) -- (1.16,2) -- (20,2) -- (20,3) -- (21.16,3) -- (21.16,2) -- (23,2);
\draw[thick,ambar] (0,0) -- (0,1) -- (1.4,1) -- (1.4,0) -- (20,0) -- (20,1) -- (21.4,1) -- (21.4,0) -- (23,0);
\draw[<->] (0,-0.6) -- node[below,lab]{20 ms = 1\,000\,000 ciclos} (20,-0.6);
\node[lab,anchor=west] at (1.5,3.3) {1.16 ms (izquierda)};
\node[lab,anchor=west] at (1.6,1.3) {1.40 ms};
\end{tikzpicture}
\end{center}
\end{resp}

\preg{M}{¿Cómo genera \texttt{servo\_controller} seis PWM distintos con un solo contador?}
\begin{resp}
Un contador \cod{pwm\_count} va de 0 a 999\,999 (20 ms) y se reinicia. Cada canal
compara: \cod{servo\_pwm[i] = (pwm\_count < ancho\_i)}, donde \cod{ancho\_i} es el de
``arriba'' si su bit vale 1 y el de ``abajo'' si vale 0. Como todos tienen el mismo
periodo, un contador compartido basta y ahorra seis contadores de 20 bits. El
\cod{generate for} fabrica seis copias del comparador.
\end{resp}

\preg{A}{¿Por qué los servos izquierdos y derechos tienen anchos distintos?}
\begin{resp}
Están montados en espejo, enfrentados. Para que ambos ``levanten'' el punto, deben
girar en sentidos opuestos: la izquierda acorta el pulso (70\,000$\to$58\,000, de
1.40 a 1.16 ms) y la derecha lo alarga (80\,000$\to$100\,000, de 1.60 a 2.00 ms).
Los valores salieron de una calibración física y se fijan en \cod{braille\_top}.
\end{resp}

\preg{A}{¿Por qué el patrón se copia en \texttt{pattern\_active} solo al inicio de cada periodo?}
\begin{resp}
Si \cod{cell\_pattern} cambiara a mitad de un pulso, ese pulso podría salir con una
duración intermedia (ni la de arriba ni la de abajo), provocando un tirón del
servo. Copiándolo solo cuando \cod{pwm\_count} se reinicia, cada pulso sale completo
con un solo valor. El costo es hasta 20 ms de retraso, imperceptible.
\end{resp}

\preg{M}{¿Por qué el contador PWM es de 20 bits? ¿Cuánto equivale un grado?}
\begin{resp}
Debe llegar a 999\,999: $2^{20} = 1\,048\,576$. A 50 MHz, 1 ms = 50\,000 cuentas; si
el rango de 1 a 2 ms cubre $\sim$180°, un grado son $\approx 50\,000/180 \approx 278$
cuentas.
\end{resp}

\preg{M}{¿Dónde se decide realmente el movimiento de los servos?}
\begin{resp}
En tres lugares: \textbf{qué} puntos suben lo decide el controller
(\cod{cell\_pattern}); \textbf{cuánto} giran lo fijan los parámetros
\cod{PULSE\_*} del top; \textbf{cómo} se genera la señal lo hace
\cod{servo\_controller}. El top solo conecta.
\end{resp}

% =====================================================================
\section{El controlador: \texttt{braille\_controller}}

\preg{M}{Dibuja la máquina de estados del controller.}
\begin{resp}
\begin{center}
\begin{tikzpicture}[->,auto]
\node[sc] (I) at (0,0) {ST\_IDLE};
\node[sc] (C) at (3.6,2.2) {ST\_CMD};
\node[sc] (L) at (7.6,2.2) {ST\_LDIG2};
\node[sc] (S) at (3.6,0) {ST\_SHOW};
\node[sc] (T) at (3.6,-2.2) {ST\_TARGET};
\node[sc] (N) at (11.4,0) {ST\_SEND};
\node[sc] (G) at (11.4,-2.4) {ST\_GAP};
\draw (I) -- node[lab,above left]{line\_ready} (C);
\draw (I) -- node[lab,above]{NEXT (lect.)} (S);
\draw (I) -- node[lab,below left]{need\_target} (T);
\draw (C) -- node[lab,above]{L de 2 dígitos} (L);
\draw (C) -- node[lab,left]{R válido} (S);
\draw (C.south east) to[out=-15,in=165] node[lab,above,pos=0.6]{resto} (N.north west);
\draw (L) -- (N);
\draw (S) -- node[lab,above]{D c} (N);
\draw (T.east) to[out=0,in=-150] node[lab,below]{T c} (N.south west);
\draw (N) edge[loop right] node[lab]{busy} (N);
\draw (N) -- node[lab,right]{msg\_start} (G);
\draw (G.south) -- ++(0,-0.7) -| node[lab,below,pos=0.25]{1 ciclo} (I.south);
\draw (I.north) -- ++(0,3.4) -| node[lab,above,pos=0.25]{modo, acorde, NEXT, F: mensaje directo} (N.north);
\end{tikzpicture}
\end{center}
Siete estados codificados en 3 bits. Casi todos los caminos terminan en
ST\_SEND$\to$ST\_GAP$\to$ST\_IDLE, porque casi todo genera un mensaje para MATLAB.
El comando \cod{N} vuelve directo a IDLE marcando \cod{next\_pend}.
\end{resp}

\preg{A}{¿En qué orden atiende los eventos ST\_IDLE y por qué ese orden?}
\begin{resp}
\begin{tabularx}{\linewidth}{c l X}
\toprule
\textbf{Prioridad} & \textbf{Evento} & \textbf{Acción} \\
\midrule
1 & Línea de la PC (\cod{line\_ready}) & ack, copia y ST\_CMD \\
2 & Cambio de switches & nuevo modo, limpia, \cod{M} + modo \\
3 & Fin de sesión (\cod{need\_final}) & \cod{F} + marcador \\
4 & Letra nueva (\cod{need\_target}) & LFSR, ST\_TARGET \\
5 & Acorde (\cod{chord\_pend}) & \cod{K}, o compara y \cod{O}/\cod{X} \\
6 & NEXT (\cod{next\_pend}) & según modo \\
\bottomrule
\end{tabularx}
La PC primero, porque mientras \cod{line\_ready = 1} el receptor bloquea bytes nuevos:
liberarlo pronto evita perder datos. El cambio de modo antes que los acordes, para
no interpretar un acorde con las reglas del modo viejo. El fin de sesión antes que
pedir otra letra. Se atiende \textbf{un evento por vuelta}; los demás esperan.
\end{resp}

\preg{A}{¿Cómo se garantiza que no se pierda un acorde o un NEXT mientras el controller está ocupado?}
\begin{resp}
Fuera del \cod{case}, en cada ciclo:
\begin{lstlisting}
if (chord_valid) begin chord_pend <= 1; chord_reg <= chord; end
if (next_pulse)  next_pend  <= 1;
\end{lstlisting}
El evento se guarda como \emph{pendiente} aunque la FSM esté en otro estado, y se
atiende al volver a ST\_IDLE. Además el acorde se atiende con
\cod{chord\_pend \&\& !chord\_valid}, es decir, al menos un ciclo después de capturarlo,
cuando \cod{chord\_reg} ya tiene el valor estable.
\end{resp}

\preg{M}{¿Qué hace el controller al cambiar los switches?}
\begin{resp}
Compara \cod{sw\_mode} con \cod{sw\_last}; si cambió: actualiza \cod{mode}, pone la
celda en vacío, borra la palabra (\cod{word\_len = 0}), limpia todas las banderas de
práctica y pendientes, y envía \cod{M} + modo. Si el nuevo modo es práctica, además
borra \cod{hits} y \cod{misses} y pide la primera letra (\cod{need\_target = 1}).
\end{resp}

\preg{M}{Explica el modo lectura de principio a fin.}
\begin{resp}
\cod{R<texto>} $\to$ ST\_CMD valida: modo lectura (si no, \cod{!M}), sin overflow (si
no, \cod{!L} y borra la palabra), largo $>$ 0 (si no, \cod{!U}). Guarda
\cod{word\_len} y va a ST\_SHOW con \cod{char\_idx = 0}: lee la casilla, la convierte con
\cod{ascii\_to\_braille}, pone \cod{cell\_pattern} y envía \cod{D} + carácter. Cada
NEXT incrementa \cod{char\_idx} y repite ST\_SHOW; tras la última, envía \cod{E},
pone \cod{word\_len = 0} y baja los servos.
\end{resp}

\preg{M}{Explica el modo escritura.}
\begin{resp}
Cada acorde se convierte con \cod{braille\_to\_ascii(chord\_reg)} y se envía
\cod{K} + letra (\cod{?} si el patrón no existe); también se guarda en
\cod{last\_chord} para los LEDs. NEXT envía \cod{K} + espacio. Los servos quedan
abajo.
\end{resp}

\preg{A}{Explica el modo práctica completo, con su flujo.}
\begin{resp}
\begin{center}
\begin{tikzpicture}[node distance=0.5cm and 0.8cm,->]
\node[act] (a) {entrar a práctica\\hits = misses = 0};
\node[act,below=of a] (b) {LFSR elige letra\\envía T + letra};
\node[act,below=of b] (c) {espera acorde\\(pista si SW2 = 1)};
\node[dec,below=of c] (d) {¿acorde =\\patrón?};
\node[act,left=1.2cm of d] (e) {hits + 1\\envía O};
\node[act,right=1.2cm of d] (f) {misses + 1\\envía X};
\node[act,below=1.3cm of d] (g) {muestra respuesta 2 s};
\node[dec,below=of g] (h) {¿último\\intento?};
\node[act,right=1.2cm of h] (i) {envía F\\espera NEXT};
\draw (a) -- (b); \draw (b) -- (c); \draw (c) -- (d);
\draw (d) -- node[lab,above]{sí} (e); \draw (d) -- node[lab,above]{no} (f);
\draw (e) |- (g); \draw (f) |- (g); \draw (g) -- (h);
\draw (h) -- node[lab,above]{sí} (i);
\draw (h.west) -- ++(-2.4,0) |- node[lab,above,pos=0.75]{no} (b.west);
\draw (i.north) -- ++(0,0.5) -| node[lab,right,pos=0.3]{NEXT: nueva sesión} ++(1.3,0) |- (a.east);
\draw[dashed] (c.east) -- ++(2.8,0) node[lab,right]{NEXT: salta\\(no cuenta)};
\end{tikzpicture}
\end{center}
\end{resp}

\preg{A}{¿Cómo se decide si un intento es acierto o error?}
\begin{resp}
\begin{lstlisting}
wire [7:0] lut_in = (mode == MODE_PRACT) ? target : buf_rdata;
ascii_to_braille A2B (.ascii_in(lut_in), .braille_out(lut_pattern), ...);
wire is_hit = (chord_reg == lut_pattern);
\end{lstlisting}
En práctica, la LUT convierte la letra objetivo en su patrón correcto y se compara
bit a bit con el acorde: deben coincidir los 6 puntos. Solo se evalúa si
\cod{target\_valid = 1} (hay una letra esperando respuesta); al responder se pone en
0, así que un segundo acorde durante los 2 s de retroalimentación se ignora.
\end{resp}

\preg{M}{¿Por qué una sola instancia de \texttt{ascii\_to\_braille} sirve para lectura y práctica?}
\begin{resp}
Porque los dos modos son excluyentes. El multiplexor \cod{lut\_in} elige la entrada
según el modo: \cod{buf\_rdata} en lectura, \cod{target} en práctica. Se ahorra una
segunda tabla.
\end{resp}

\preg{A}{Explica cómo funciona el LFSR, con diagrama.}
\begin{resp}
\begin{center}
\begin{tikzpicture}[scale=0.62]
\foreach \i in {0,...,15}{
  \pgfmathtruncatemacro{\b}{15-\i}
  \draw[thick,fill=white] (\i,0) rectangle (\i+1,1);
  \node[font=\tiny\ttfamily] at (\i+0.5,0.5) {\b};}
\foreach \i in {0,2,3,5}{\fill[verdesuave] (\i,0) rectangle (\i+1,1); \draw[thick] (\i,0) rectangle (\i+1,1);
  \pgfmathtruncatemacro{\b}{15-\i}\node[font=\tiny\ttfamily] at (\i+0.5,0.5) {\b};
  \draw[thick,verde] (\i+0.5,1) -- (\i+0.5,2);}
\draw[thick,verde] (0.5,2) -- (17.5,2);
\node[circle,draw=verde,thick,fill=white,inner sep=1pt,font=\small] (x) at (17.5,1.2) {$\oplus$};
\draw[->,thick,verde] (17.5,2) -- (x);
\draw[->,thick,verde] (x.south) |- (16,0.5);
\draw[->,thick,gris] (15.5,-0.4) -- node[below,lab]{recorre a la izquierda en cada ciclo} (0.5,-0.4);
\end{tikzpicture}
\end{center}
\begin{lstlisting}
lfsr <= {lfsr[14:0], lfsr[15] ^ lfsr[13] ^ lfsr[12] ^ lfsr[10]};
\end{lstlisting}
Registro de 16 bits que en cada ciclo se recorre una posición y mete por la derecha
el XOR de los bits 15, 13, 12 y 10. Esas \emph{taps} corresponden a un polinomio
primitivo ($x^{16}+x^{14}+x^{13}+x^{11}+1$), por lo que recorre los $2^{16}-1 =
65\,535$ estados no nulos antes de repetirse (\emph{maximal length}). La semilla
\cod{0xACE1} no puede ser 0: con todo en cero el XOR siempre da 0 y se bloquea.
\end{resp}

\preg{A}{Si la secuencia del LFSR es fija, ¿por qué las letras parecen aleatorias?}
\begin{resp}
El LFSR avanza 50 millones de veces por segundo y la letra se toma en el instante
en que hace falta (cuando el usuario termina un intento y pasan los 2 s, o presiona
NEXT). Ese instante depende del tiempo humano, que nadie controla con precisión de
nanosegundos. Es pseudoaleatorio, pero impredecible en la práctica.
\end{resp}
\begin{trampa}
Después de un reset la secuencia arranca igual, pero la letra elegida depende de
cuánto tiempo pasa hasta que se necesita, así que no se repite la misma sesión.
\end{trampa}

\preg{A}{¿Cómo se convierte el valor del LFSR en una letra? ¿La distribución es uniforme?}
\begin{resp}
\begin{lstlisting}
wire [4:0] r5       = lfsr[4:0];                     // 0..31
wire [4:0] rnd_idx  = (r5 >= 26) ? (r5 - 26) : r5;   // 0..25
wire [7:0] rnd_char = 8'h61 + rnd_idx;               // 'a'..'z'
wire [7:0] new_char = (rnd_char != target) ? rnd_char :
                      (rnd_idx == 25) ? 8'h61 : rnd_char + 1;  // no repetir
\end{lstlisting}
Las letras son consecutivas en ASCII, así que basta sumar el índice a \cod{'a'}.
\textbf{No es uniforme}: los valores 26--31 se doblan a 0--5, así que \cod{a}--\cod{f}
tienen probabilidad 2/32 y las demás 1/32. Para práctica es aceptable; una solución
sería descartar y volver a muestrear, o usar más bits y módulo.
\end{resp}

\preg{M}{¿Por qué \texttt{0 = sin límite} en \texttt{L}?}
\begin{resp}
\begin{lstlisting}
last_try <= (rounds_cfg != 0) && (total_new >= rounds_cfg);
\end{lstlisting}
Una sesión de 0 letras no tiene sentido, así que ese valor se reutiliza como caso
especial sin necesitar otro comando. Con 0, \cod{rounds\_cfg != 0} es falso, nunca
hay último intento ni \cod{F}. Sin esa condición, \cod{total\_new >= 0} sería verdadero
desde el primer intento.
\end{resp}

\preg{A}{¿Cómo se interpreta el número en \texttt{L10}? ¿Por qué \texttt{digit\_val = buf\_rdata[3:0]}?}
\begin{resp}
En ST\_CMD se valida que haya 1 o 2 caracteres y que el primero sea dígito
(\cod{'0'} $\le$ c $\le$ \cod{'9'}). Con uno solo, \cod{rounds\_cfg = dígito}. Con dos,
guarda el primero en \cod{dig0}, pasa a ST\_LDIG2 (el buffer tiene lectura
combinacional, pero la dirección cambia a la casilla 1) y calcula
\cod{dig0*10 + dígito}. Responde con el valor en hex (\cod{L10} $\to$ \cod{L0A}).\par
Los dígitos ASCII son \cod{0x30}--\cod{0x39}: el nibble bajo es exactamente el valor
del dígito, así que \cod{buf\_rdata[3:0]} lo extrae sin restar.
\end{resp}

\preg{M}{¿Cómo funciona la pista (SW2)?}
\begin{resp}
\begin{lstlisting}
if (fb_wait)                      cell_pattern <= lut_pattern;  // respuesta 2 s
else if (target_valid && sw_hint) cell_pattern <= lut_pattern;  // pista
else                              cell_pattern <= 6'b000000;
\end{lstlisting}
Se evalúa en cada ciclo en modo práctica, así que mover SW2 aplica de inmediato. La
pista no afecta la evaluación: el acierto se cuenta igual.
\end{resp}

\preg{M}{¿Cómo funciona el temporizador de 2 segundos?}
\begin{resp}
Tras un intento, \cod{fb\_wait = 1} y \cod{fb\_cnt} cuenta hasta
\cod{FEEDBACK\_DELAY - 1} = 99\,999\,999 (2 s a 50 MHz; 27 bits, $2^{27}\approx 134$
millones). Al llegar, baja \cod{fb\_wait} y levanta \cod{need\_final} (si era el
último intento) o \cod{need\_target} (siguiente letra). Corre fuera del \cod{case}, en
paralelo con la FSM.
\end{resp}

\preg{A}{¿Qué pasa si se presiona NEXT en cada momento de práctica?}
\begin{resp}
Esperando respuesta: salta la letra sin contarla (\cod{target\_valid = 0},
\cod{need\_target = 1}). Durante los 2 s de retroalimentación: se ignora. Con la
sesión terminada (\cod{sess\_done}): empieza una sesión nueva (borra el marcador,
envía \cod{MP} y pide letra).
\end{resp}

\preg{A}{¿Cuál es el rango de \texttt{hits} y \texttt{misses} y qué pasa al desbordar?}
\begin{resp}
8 bits: 0 a 255. Al pasar de 255 vuelven a 0 (aritmética módulo 256) sin aviso. Con
sesiones de hasta 99 letras no ocurre, pero con \cod{L0} (sin límite) es posible
tras 256 intentos. Mejora: saturar en 255 o ampliar a 16 bits (4 dígitos hex).
\end{resp}

% =====================================================================
\section{Integración: \texttt{braille\_top}}

\preg{B}{¿Qué hace \texttt{braille\_top}? ¿Tiene lógica propia?}
\begin{resp}
Nombra los pines físicos, fija todos los parámetros e instancia y conecta los
módulos. Su única lógica propia es \cod{rst = \textasciitilde KEY[0]}, el sincronizador
de dos flip-flops de los switches y los \cod{assign} de LEDs y puntos decimales.
\end{resp}

\preg{M}{¿Por qué los switches se sincronizan pero no tienen antirrebote?}
\begin{resp}
Se sincronizan porque son asíncronos (riesgo de metaestabilidad). No se les quita
el rebote porque se mueven raramente y un rebote solo provocaría, como mucho,
cambios de modo muy breves: el controller enviaría algún \cod{M} extra y terminaría
en el modo correcto. Es una simplificación consciente; un antirrebote los haría más
robustos.
\end{resp}

\preg{M}{¿Qué muestran los LEDs y displays y para qué sirven?}
\begin{resp}
\cod{LEDR[5:0]}: el patrón de la celda; en escritura, los botones presionados en vivo
o el último acorde. \cod{LEDR[7:6]} y HEX5: modo. \cod{LEDR[8]}: NEXT presionado.
\cod{LEDR[9]}: transmitiendo (\cod{msg\_busy}). HEX3--HEX2: aciertos; HEX1--HEX0:
errores; HEX4 apagado. Sirven para depurar sin MATLAB: si un LED se enciende pero su
servo no se mueve, el problema es mecánico o de cableado, no lógico.
\end{resp}

\preg{M}{¿Cómo se agrupan los botones en \texttt{\{KEY[1], DOT\_N\}}?}
\begin{resp}
La concatenación forma un vector de 7 bits: bit 6 = NEXT, bits 5..0 = puntos 6..1.
\cod{btn\_pressed[5:0]} va a \cod{chord\_capture} y \cod{btn\_pulse[6]} al controller.
Un solo \cod{button\_conditioner} parametrizado (\cod{N\_BUTTONS = 7}) atiende a todos.
\end{resp}

% =====================================================================
\section{La interfaz MATLAB}

\preg{B}{¿Qué hace MATLAB y qué envía?}
\begin{resp}
Abre el puerto serial, escucha, muestra en pantalla, habla y grafica. Envía solo
tres cosas: \cod{Q} al conectar, \cod{R} + palabra en lectura y \cod{L} + número en
práctica.
\end{resp}

\preg{M}{¿Cómo detecta MATLAB el fin de un mensaje y cómo maneja mensajes juntos o partidos?}
\begin{resp}
\cod{configureTerminator(s,"LF")} y \cod{configureCallback(s,"terminator",@alRecibir)}
ejecutan \cod{alRecibir} al llegar un \cod{\textbackslash n}. Esa función acumula los
bytes en \cod{rxbuf} y, en un \cod{while}, corta un mensaje por cada byte 10. Los
mensajes completos se procesan; un pedazo incompleto se queda en el buffer hasta que
llegue su terminador. Es el equivalente de \cod{line\_rx} en la PC.
\end{resp}

\preg{M}{¿Por qué MATLAB envía \texttt{Q} al conectarse?}
\begin{resp}
Porque el modo se elige con los switches y MATLAB no puede verlos. Al conectarse no
sabe en qué modo está la tarjeta; \cod{Q} pide a la FPGA que responda \cod{M} + modo.
\end{resp}

\preg{M}{¿Dónde se calcula la precisión y por qué ahí?}
\begin{resp}
En MATLAB: \cod{acc = 100*h/max(1,h+e)}. La FPGA cuenta (es quien sabe si acertaste)
y envía los totales; dividir con decimales en hardware sería costoso y solo sirve
para mostrar. \cod{max(1,\ldots)} evita dividir entre cero.
\end{resp}

\preg{A}{¿Por qué MATLAB no lleva su propia cuenta de aciertos?}
\begin{resp}
Para no desincronizarse: si se perdiera un mensaje, o MATLAB se conectara a mitad de
una sesión, una cuenta propia sería incorrecta. Leer siempre los totales de la FPGA
garantiza que la pantalla coincide con los displays.
\end{resp}

\preg{M}{¿Cómo funciona la voz?}
\begin{resp}
Con el sintetizador de Windows a través de .NET: \cod{NET.addAssembly('System.Speech')}
y \cod{SpeechSynthesizer}. \cod{SpeakAsync} habla en segundo plano para no congelar la
ventana; \cod{SpeakAsyncCancelAll} calla la frase anterior para que no se acumulen.
Todo está en \cod{try/catch}: sin Windows, el programa funciona sin voz.
\end{resp}

\preg{M}{¿Cómo cambia de panel la ventana?}
\begin{resp}
Los cuatro paneles están en la misma celda de la cuadrícula; \cod{mostrarModo} deja
visible solo el del modo actual con \cod{pX.Visible = (m == 'X')}.
\end{resp}

\preg{A}{¿Por qué las funciones de MATLAB comparten variables sin recibirlas como argumento?}
\begin{resp}
Son funciones anidadas dentro de \cod{braille\_gui}; en MATLAB comparten el espacio de
trabajo de la función que las contiene. Por eso \cod{s}, \cod{rxbuf} o \cod{histAcc}
se ven desde todos los callbacks.
\end{resp}

\preg{M}{¿Qué protección hay si la palabra tiene más de 16 caracteres?}
\begin{resp}
Doble: MATLAB la rechaza antes de enviar (\cod{numel(w) > MAX\_LEN} $\to$
\cod{uialert}); si llegara desde otro programa, \cod{line\_rx} marca overflow y el
controller responde \cod{!L}.
\end{resp}

% =====================================================================
\section{Análisis crítico, límites y mejoras}

\preg{A}{¿Cuáles son las principales limitaciones del diseño?}
\begin{resp}
\begin{itemize}
  \item Sin FIFO en recepción: depende de que la PC espere cada respuesta.
  \item Sin paridad ni detección de errores más allá del STOP; los errores de trama
        se descartan sin avisar.
  \item Contadores de 8 bits que dan la vuelta en 256.
  \item Distribución del LFSR sesgada (a--f con el doble de probabilidad).
  \item Reset asíncrono sin sincronizar su liberación; switches sin antirrebote.
  \item Sin signo de mayúscula ni números en Braille (solo grado 1 básico).
  \item Voz dependiente de Windows.
\end{itemize}
\end{resp}

\preg{A}{Si tuvieras que mejorar el proyecto, ¿qué harías primero y por qué?}
\begin{resp}
(1) Un sincronizador de reset y antirrebote para switches: robustez con muy poco
costo. (2) Una FIFO pequeña en recepción para tolerar comandos seguidos. (3)
Saturación o 16 bits en contadores. (4) Muestreo sin sesgo en el LFSR. (5) Soporte de
números y mayúsculas en Braille. (6) Un testbench automatizado por módulo.
\end{resp}

\preg{A}{¿Cómo verificaste el diseño? ¿Cómo lo verificarías con más rigor?}
\begin{resp}
Pruebas en hardware con los LEDs y displays como observadores, y comandos de prueba
(\cod{S}, \cod{C}, \cod{N}) desde un monitor serial. Con más rigor: testbenches en
simulación (ModelSim/Questa) por módulo, reduciendo parámetros (por ejemplo
\cod{CLKS\_PER\_BIT = 4}, \cod{DEBOUNCE\_COUNT = 8}) para simular rápido; un modelo de
UART en el testbench que envíe tramas y compruebe las respuestas; casos límite
(17 caracteres, \cod{L100}, comandos en minúscula, STOP erróneo) y SignalTap en la
tarjeta.
\end{resp}

\preg{A}{¿Qué pasaría si cambias la velocidad a 115200 baudios? ¿Qué tendrías que modificar?}
\begin{resp}
\cod{CLKS\_PER\_BIT} = $50\,000\,000/115\,200 = 434.03 \to 434$ (error 0.007\%), en el
top, y \cod{BAUD} en MATLAB. Nada más: el resto del diseño no depende del tiempo de
bit. El contador de 13 bits sobra. Un byte tardaría 87\,µs.
\end{resp}

\preg{A}{¿Qué pasa si se presionan botones durante la transmisión de un mensaje?}
\begin{resp}
No se pierden: \cod{chord\_valid} y \cod{next\_pulse} se guardan como pendientes fuera
del \cod{case}, y se atienden cuando la FSM vuelve a ST\_IDLE. La transmisión no se
interrumpe porque \cod{msg\_tx} tiene su copia del mensaje.
\end{resp}

\preg{A}{¿Qué pasa si se cambia de modo a mitad de una palabra o de una sesión?}
\begin{resp}
El cambio de switches tiene prioridad 2: se limpia la palabra (\cod{word\_len = 0}),
la celda, las banderas de práctica y los pendientes, y se envía el nuevo modo. Al
volver a práctica empieza una sesión nueva con el marcador en cero.
\end{resp}

\preg{A}{¿Hay algún riesgo si llegan dos acordes muy seguidos?}
\begin{resp}
\cod{chord\_capture} necesita que se suelten todos los botones entre acordes y el
antirrebote impone 10 ms por transición, así que dos acordes están separados por
decenas de milisegundos, mientras que el controller atiende un acorde en pocos
ciclos más el envío. Si llegara un segundo acorde con \cod{chord\_pend} todavía en
1, sobrescribiría \cod{chord\_reg} y solo se procesaría el último: hay una sola
posición de pendiente, no una cola. En la práctica, con tiempos humanos, no ocurre.
\end{resp}

\preg{A}{¿Por qué usar una FSM en el controller y no solo lógica combinacional?}
\begin{resp}
Porque las acciones son secuenciales y dependen de lo que pasó antes: leer el
buffer, convertir, esperar a que el transmisor esté libre, enviar, volver a esperar.
Una FSM hace explícitos esos pasos, asegura que solo se atienda un evento a la vez y
que cada mensaje se complete antes del siguiente.
\end{resp}

\preg{A}{Justifica el uso de 9600 baudios.}
\begin{resp}
Es una velocidad estándar que aceptan MATLAB, el FT232RL y el HC-05 sin
configuración. El tráfico es mínimo (mensajes de 2 a 18 bytes al ritmo de una
persona): a 960 bytes/s el mensaje más largo tarda 18.8 ms. Velocidades menores dan
más margen frente a errores de reloj y ruido en cables.
\end{resp}

% =====================================================================
\section{Ejercicios de traza (respuesta paso a paso)}

\preg{A}{En práctica llevas 2 aciertos y 1 error. La letra es \texttt{m} y tecleas los puntos 1-3-4. Traza todo.}
\begin{resp}
\begin{enumerate}
  \item \cod{button\_conditioner}: tras 10 ms estables, \cod{btn\_pressed} muestra los
        puntos 1, 3 y 4.
  \item \cod{chord\_capture}: acumula \cod{001101}; al soltar todos, \cod{chord = 001101}
        y pulso \cod{chord\_valid}.
  \item Controller: \cod{chord\_pend = 1}, \cod{chord\_reg = 001101}. En ST\_IDLE (prioridad
        5), modo práctica y \cod{target\_valid = 1}.
  \item \cod{lut\_pattern} = \cod{ascii\_to\_braille('m')} = \cod{001101}; \cod{is\_hit = 1}.
  \item \cod{hits = 3}, \cod{target\_valid = 0}, \cod{fb\_wait = 1},
        \cod{msg\_data = "Om0301"}, \cod{msg\_len = 6}, ST\_SEND.
  \item ST\_SEND$\to$\cod{msg\_start}$\to$ST\_GAP$\to$ST\_IDLE. \cod{msg\_tx} envía 7 bytes
        (\cod{O m 0 3 0 1 \textbackslash n}, 7.3 ms).
  \item HEX3--HEX2 = \cod{03}; la celda muestra la \cod{m} durante 2 s.
  \item MATLAB: \cod{h=3}, \cod{e=1}, precisión 75\%, letra en verde, ``Correcto''.
  \item A los 2 s: \cod{need\_target = 1} $\to$ nueva letra $\to$ \cod{T} + letra.
\end{enumerate}
\end{resp}

\preg{A}{En modo escritura envías por error \texttt{Rhola} desde MATLAB. ¿Qué responde y por qué?}
\begin{resp}
\cod{line\_rx} guarda la palabra en el buffer (el comando es \cod{R}) y levanta
\cod{line\_ready}. El controller da \cod{ack} y en ST\_CMD ve \cod{mode != MODE\_READ}:
responde \cod{!M}. La palabra queda en el buffer pero no se usa (\cod{word\_len} no
cambia). MATLAB muestra ``Pon la tarjeta en modo lectura''.
\end{resp}

\preg{A}{Envías \texttt{R} seguido de 17 letras. ¿Qué pasa en cada módulo?}
\begin{resp}
\cod{line\_rx}: guarda las primeras 16 en las casillas 0--15 (\cod{len} llega a 16);
la 17.ª activa \cod{overflow}; el \cod{\textbackslash n} levanta \cod{line\_ready}. El
controller, en ST\_CMD con \cod{ovf\_r = 1}: \cod{word\_len = 0}, celda vacía, responde
\cod{!L}. (MATLAB normalmente lo impide antes de enviar.)
\end{resp}

\preg{A}{Envías \texttt{L5}, luego \texttt{L20}, luego \texttt{L100}. ¿Qué responde cada uno?}
\begin{resp}
\cod{L5} $\to$ \cod{L05} (\cod{rounds\_cfg = 5}). \cod{L20} $\to$ pasa por ST\_LDIG2,
$2\cdot10+0 = 20 = $ \cod{0x14} $\to$ \cod{L14}. \cod{L100} $\to$ \cod{len = 3 > 2}
$\to$ \cod{!U} y \cod{rounds\_cfg} no cambia.
\end{resp}

\preg{A}{Llega por el cable un byte con ruido que baja la línea solo 10 µs. ¿Qué hace \texttt{uart\_rx}?}
\begin{resp}
El sincronizador lo deja pasar (dura 500 ciclos) y la FSM entra a START. A los 52 µs
(medio bit) la línea ya volvió a 1: se considera falso START y regresa a IDLE sin
entregar nada.
\end{resp}

\preg{A}{Estás en lectura con ``hi'' cargada y presionas NEXT tres veces. ¿Qué sale por la UART?}
\begin{resp}
Al cargar: \cod{Dh}. 1.er NEXT: \cod{char\_idx = 1} $\to$ \cod{Di}. 2.º NEXT:
\cod{char\_idx + 1 = 2} no es menor que \cod{word\_len = 2} $\to$ \cod{E}, celda vacía,
\cod{word\_len = 0}. 3.er NEXT: no hay palabra (\cod{word\_len = 0}), no hace nada ni
responde.
\end{resp}

% =====================================================================
\section{Preguntas relámpago}

\begin{longtable}{p{8.2cm} p{7cm}}
\toprule
\textbf{Pregunta} & \textbf{Respuesta corta} \\
\midrule
\endhead
Periodo del reloj de la DE10-Lite & 20 ns (50 MHz) \\
Ciclos por bit a 9600 & 5208 \\
Tiempo por carácter UART & 1.04 ms (10 bits) \\
Línea más larga que acepta la FPGA & R + 16 + \textbackslash n = 18 bytes \\
Mensaje más largo que envía & 6 caracteres + \textbackslash n = 7 bytes \\
Ancho de \cod{msg\_data} y por qué & 48 bits = 6 bytes, mensaje O/X \\
Ancho de \cod{msg\_len} & 3 bits (hasta 7) \\
Ancho de \cod{len} en \cod{line\_rx} & 5 bits (debe valer 16) \\
Bits del contador PWM & 20 (hasta 999\,999) \\
Bits del contador de antirrebote & 19 (hasta 500\,000) \\
Bits del temporizador de 2 s & 27 \\
Estados del controller & 7 (3 bits) \\
Estados de \cod{uart\_rx} / \cod{msg\_tx} & 4 / 3 \\
Semilla del LFSR & \cod{0xACE1} (no puede ser 0) \\
Periodo del LFSR & 65\,535 \\
Patrón de la \cod{h} & \cod{010011} (puntos 1, 2, 5) \\
Patrón de la ñ & \cod{111011} (puntos 1, 2, 4, 5, 6) \\
¿Qué hace \cod{blank} en \cod{hex7seg}? & Apaga el display (HEX4) \\
¿Display activo en alto o en bajo? & En bajo \\
Valor de \cod{rounds\_cfg} tras reset & 10 \\
Respuesta a \cod{C} & \cod{S0000} \\
Comando que no tiene respuesta propia & \cod{N} \\
¿Quién genera el \textbackslash n de las respuestas? & \cod{msg\_tx} \\
¿Quién guarda la palabra? & \cod{char\_buffer} \\
¿Quién decide qué puntos suben? & \cod{braille\_controller} \\
¿Dónde se calibran los servos? & Parámetros \cod{PULSE\_*} de \cod{braille\_top} \\
\bottomrule
\end{longtable}

\end{document}
